% !TEX root = ../../hw07.tex
% Source: Lucas's supplied typed solutions, October 30, 2026, 4.1--4.3.
% The completion for the alternative w=sv in 4.3 was explicitly authorized
% by Lucas on October 30, including the s=0 case.
\begin{proof}[Solution]
  \textbf{4.1.} A quadratic \(f(t)=at^2+bt+c\) with \(a>0\) is
  an upward-opening parabola and has a unique minimum. By bilinearity
  and symmetry of the real inner product,
  \[
    \begin{aligned}
      f(t)&=|v-tw|^2=(v-tw,v-tw)\\
        &=(v,v)-t(v,w)-t(w,v)+t^2(w,w)\\
        &=(w,w)t^2-2(v,w)t+(v,v).
    \end{aligned}
  \]
  Since \(w\ne0\), positive definiteness gives \(a=(w,w)>0\).
  Thus \(f\) has a unique minimum.

  \textbf{4.2.} Complete the square in \(f(t)=at^2+bt+c\). Factoring
  out \(a\) from the first two terms gives
  \[
    f(t)=a\left(t^2+\frac{b}{a}t\right)+c.
  \]
  Since \((t-d)^2=t^2-2dt+d^2\), the required shift satisfies
  \(-2d=b/a\), or \(d=-b/(2a)\). Adding and subtracting
  \((b/(2a))^2\) gives
  \[
    \begin{aligned}
      f(t)&=a\left(t^2+\frac{b}{a}t
        +\left(\frac{b}{2a}\right)^2
        -\left(\frac{b}{2a}\right)^2\right)+c\\
        &=a\left(\left(t+\frac{b}{2a}\right)^2
          -\left(\frac{b}{2a}\right)^2\right)+c\\
        &=a\left(t+\frac{b}{2a}\right)^2
          +\left(c-\frac{b^2}{4a}\right).
    \end{aligned}
  \]
  Because \(a>0\), the squared term has its unique minimum at zero.
  Consequently,
  \[
    t_{\min}=-\frac{b}{2a},\qquad
    f_{\min}=c-\frac{b^2}{4a}.
  \]
  Substituting \(a=(w,w)\), \(b=-2(v,w)\), and \(c=(v,v)\) yields
  \[
    t_{\min}=\frac{(v,w)}{(w,w)},\qquad
    f_{\min}=(v,v)-\frac{(v,w)^2}{(w,w)}.
  \]

  \textbf{4.3.} First suppose \(w\ne0\). The preceding calculation
  also applies when \(v=0\). Positive definiteness gives
  \[
    f(t)=|v-tw|^2\ge0\qquad\text{for every }t\in\R.
  \]
  Thus its global minimum satisfies
  \[
    f_{\min}=(v,v)-\frac{(v,w)^2}{(w,w)}\ge0.
  \]
  Since \((w,w)>0\), multiplying by \((w,w)\) gives
  \[
    (v,w)^2\le(v,v)(w,w).
  \]
  If equality holds, then \(f_{\min}=0\), so
  \[
    |v-t_{\min}w|^2=0.
  \]
  Positive definiteness implies \(v-t_{\min}w=0\). Hence
  \(v=sw\) with \(s=t_{\min}\).

  Conversely, if \(v=sw\) for some \(s\in\R\), then
  \[
    f(s)=|v-sw|^2=0.
  \]
  Since \(f(t)\ge0\) for every \(t\), its minimum is zero. Therefore
  \[
    (v,v)-\frac{(v,w)^2}{(w,w)}=0,
  \]
  which gives equality in the Cauchy--Schwarz inequality.
  If instead \(w=sv\), our assumption \(w\ne0\) forces
  \(s\ne0\). Then \(v=s^{-1}w\), so the same argument applies.

  Finally, if \(w=0\), bilinearity gives \((v,0)=0\) and
  \((0,0)=0\), so
  \[
    (v,w)^2=0=(v,v)(w,w).
  \]
  Equality holds, and the required alternative \(w=sv\) is satisfied
  by \(s=0\). This includes the case \(w=sv\) with \(s=0\) and
  completes the inequality and its equality characterization for all
  \(v,w\in V\).
\end{proof}
